% B proof; root and D full independent literal audits PASS.
\subsection{Bouquet graphs in the right-zero chart construction}

\begin{proposition}[A bouquet on one side of a chart]
\label{C:bouquet-chart-rank-one-square}
Use the three corner monoids of
Theorem~\ref{C:affine-chart-family-zero}: the initial corner is
the identity-adjoined left-zero two-band, the middle corner is
trivial, and the final corner is the identity-adjoined right-zero
two-band.  Let $A,B$ be the natural graphs of the two adjacent
acts, and let $K=K(T)$ be the cubical model of the long-edge
biset.  Suppose every component of $A$ is a bouquet of loops;
isolated vertices are allowed.  The graph $B$ and the composition
may be arbitrary.  If the whole category has rational second
cohomology of dimension at most one, its degree-two square is
zero.  The same conclusion holds with $A,B$ interchanged.
\end{proposition}
\begin{proof}
Write $X$ for the actual triangle-of-spaces model with vertex
spaces points, edge spaces $A,B,K$, face space $A\times B$,
and long face map $h:A\times B\to K$.  The other two face
maps are the projections.  This is the natural model supplied
by Lemmas~\ref{C:right-zero-action-graph} and
\ref{C:two-retraction-cubical-model}.

Replace each edge space by its component set and the face by
$\pi_0A\times\pi_0B$, retaining the component map induced by
$h$.  Call the resulting two-dimensional triangle space $X_0$.
The actual component maps give $q:X\to X_0$.  The strict-face
chain filtration gives the exact cohomology sequence
\[
 H^2(X_0;\Q)\xrightarrow{q^*}H^2(X;\Q)
 \longrightarrow
 \left\{(\alpha,\beta,\lambda):
 h^*\lambda=\operatorname{pr}_A^*\alpha+
                 \operatorname{pr}_B^*\beta\right\}
 \longrightarrow0,
\]
where $\alpha\in H^1(A;\Q)$, $\beta\in H^1(B;\Q)$,
and $\lambda\in H^1(K;\Q)$.  Indeed the bottom row of the
filtration is the cellular complex of $X_0$, and the next
row is the restriction map from the three edge first-cohomology
groups to $H^1(A\times B)$.  There is no vertex cohomology in positive
degree and hence no further differential in total degree two.
If $q^*$ has nonzero image, that image spans the assumed
one-dimensional group.  Its square vanishes because $X_0$
is two-dimensional.

We may therefore assume $q^*=0$.  Every rational degree-two
class is then uniquely described by a displayed triple.
After multiplying by a nonzero integer, its three entries
are integral.  Represent them by actual maps
\[
 \alpha:A\to S^1,\qquad \beta:B\to S^1,
 \qquad \lambda:K\to S^1.
\]
Their equality in first cohomology says, on each component,
that $\lambda h$ is homotopic to the addition map
$\alpha+\beta:A\times B\to S^1$.  Thus they give a map
of triangle homotopy colimits to the universal triangle
whose three edge spaces are $S^1$, whose face is
$S^1\times S^1$, and whose long map is addition.
Write $Y$ for that universal triangle.  The same filtration
gives $H^2(Y;\Q)=\Q$ and $H_4(Y;\Q)=\Q$.
The pullback of its degree-two generator is our class:
the triples agree, and the possible difference is in
the zero image of $q^*$.

Since $K$ is two-dimensional and $A,B$ are graphs, the
degree-four filtration identifies
\[
 H_4(X;\Q)=\ker\bigl(H_2(A\times B;\Q)
                            \xrightarrow{h_*}H_2(K;\Q)\bigr).
\]
Under the map to $Y$, a class $z$ in this kernel maps to
the evaluation of $\alpha\times\beta$ on $z$.  We show
that this evaluation is zero.  This suffices by cup
naturality, without computing the square in $Y$.

Choose a cellular integral one-cocycle representing
$\lambda$ on $K$.  Write $a_i$ for the literal oriented
loop edges of $A$, based at vertices $v_i$.  Every element
of $H_2(A\times B;\Q)$ has the form
\[
 z=\sum_i a_i\otimes\gamma_i,
 \qquad \gamma_i\in Z_1(B;\Q).
\]
Use $P_0,P_1$ for the horizontal retractions and $Q_0,Q_1$
for the vertical retractions on $T$.  For an edge $b$ of
$B$, let $t_{ib}$ be the actual composition state of
$a_i$ and $b$.  Since $a_i$ is a loop,
$P_0t_{ib}=P_1t_{ib}$.  Its horizontal lower face is
$Q_0t_{ib}$ and its vertical left face is $P_0t_{ib}$;
a face in the common core is a collapsed edge and is
assigned cocycle value zero.

If this top state is horizontally degenerate, the image
of $a_i$ is a point and $\alpha(a_i)=0$.  If it is
vertically degenerate, its vertical edge has value zero.
Otherwise the cohomology identity on $A\times B$ gives
\[
 \alpha(a_i)=\lambda(Q_0t_{ib}).
\]
For the vertical cycles it gives, even when $B$ is
disconnected,
\[
 \beta(\gamma_i)
   =\sum_b\gamma_i(b)\lambda(P_0t_{ib}).
\]
Consequently evaluation of $\alpha\times\beta$ on $z$
factors through its literal cellular image in $C_2(K;\Q)$:
assign to the two-cell indexed by $t$ the number
\[
 c(t)=\lambda(Q_0t)\lambda(P_0t).
\]
Degenerate top states contribute zero as just checked.
For $z\in\ker h_*$, its cellular image is the zero
two-chain, because $K$ has no cellular three-boundaries.
Thus $\langle\alpha\times\beta,z\rangle=0$.
The induced map on fourth homology is zero, so the
pullback square vanishes over $\Q$.  Rational scaling
proves the proposition for the original class.
\end{proof}

\begin{corollary}\label{C:three-state-adjacent-chart-square}
With the same three corner monoids, if either adjacent
act has at most three states, the category cannot have
the rational cohomology ring of $\CP^2$, regardless of
the other two hom-set sizes.
\end{corollary}
\begin{proof}
A right-zero graph on at most three states is either
a forest or a disjoint union of bouquets of loops.
In the forest case its first homology is zero, so the
displayed fourth-homology kernel is zero.  In the other
case apply the proposition.
\end{proof}
